\sect{श्रवणद्वादशीव्रतम्}

\label{sec:padma-shravana-dvadashi}

\ganapatyadiDhyanam

\padmaPuranam

\dnsub{अथ सप्ततितमोऽध्यायः॥७०}

\uvacha{नारद उवाच}

\twolineshloka
{उपवासासमर्थानां सदैवामरसत्तम}
{एका या द्वादशी पुण्या तां वदस्व ममानघ}% ॥१॥

\uvacha{शिव उवाच}

\twolineshloka
{मासि भाद्रपदे शुक्ले द्वादशी श्रवणान्विता}
{सा वै सर्वप्रदा पुण्या ह्युपवासे महाफला}% ॥२॥

\twolineshloka
{सङ्गमे सरितः स्नात्वा द्वादशीं तामुपोषितः}
{अयत्नात्समवाप्नोति द्वादश द्वादशीफलम्}% ॥३॥

\twolineshloka
{बुधश्रवणसंयुक्ता या च वै द्वादशी भवेत्}
{अतीव महती तस्यां कृतं सर्वमथाक्षयम्}% ॥४॥

\twolineshloka
{द्वादशी श्रवणोपेता यदा भवति नारद}
{सङ्गमे सरितां स्नात्वा लभेद् गोदानजं फलम्}% ॥५॥

\twolineshloka
{जलपूर्णं तदा कुम्भं स्थापयित्वा विचक्षणः}
{तस्योपरि न्यसेत् पात्रं स्थापयित्वा जनार्दनम्}% ॥६॥

\twolineshloka
{ततस्तस्याग्रतो देयं नैवेद्यं घृतपाचितम्}
{सोदकांश्च नवान् कुम्भान् दद्याच्छक्त्या विचक्षणः}% ॥७॥

\twolineshloka
{एवं सम्पूज्य गोविन्दं जागरं तत्र कारयेत्}
{प्रभाते विमले स्नात्वा सम्पूज्य गरुडध्वजम्}% ॥८॥

\twolineshloka
{पुष्पधूपादि नैवेद्यैः फलैर्वस्त्रैः सुशोभनैः}
{पुष्पाञ्जलिं ततो दद्यात् मन्त्रमेनमुदीरयेत्}% ॥९॥

\twolineshloka
{नमो नमस्ते गोविन्द बुधश्रवणसंयुत}
{अघौघसंक्षयं कृत्वा सर्वसौख्यप्रदो भव}% ॥१०॥

\twolineshloka
{अन्नं तु ब्राह्मणे पूतं वेदवेदाङ्गपारगे}
{पुराणज्ञे विशेषेण विधिवत् सम्प्रदापयेत्}% ॥११॥

\twolineshloka
{अनेन विधिना चैव नद्यास्तीरे नरोत्तमः}
{सर्वं निर्वर्तयेत् सम्यग् एकचित्तरतोऽपि सन्}% ॥१२॥

\twolineshloka
{अत्राप्युदाहरन्तीमम् इतिहासं पुरातनम्}
{महत्यरण्ये यद् वृत्तं भूमिदेव शृणुष्व तत्}% ॥१३॥

\twolineshloka
{यं श्रुत्वा मानवो लोके महादुःखात् प्रमुच्यते}
{देशो दासरको नाम तस्य भागे च पश्चिमे}% ॥१४॥

\twolineshloka
{तत्र विद्वन्मरुर्देशः सर्वसत्त्वभयङ्करः}
{सुतप्तसिकता भूमिर्यत्र दुष्टा महोरगाः}% ॥१५॥

\twolineshloka
{अल्पच्छायाद्रुमाकीर्णा मृतप्राणिसमाकुला}
{शमीखदिरपालाशकरीरैः पीलुभिः सह}% ॥१६॥

\twolineshloka
{तत्र भीमा द्रुमगणाः कण्टकैरचिता दृढैः}
{जग्धप्राणजनाकीर्णा यत्र भूर्दृश्यते क्वचित्}% ॥१७॥

\twolineshloka
{तथापि जीवा जीवन्ति सर्वे कर्मनिबन्धनात्}
{नोदकं नोदकाधारा विद्वंस् तत्र बलाहकाः}% ॥१८॥

\twolineshloka
{पक्षान्तरगतैः कैश्चिच्छिशुभिस्तृषितैः समम्}
{उत्क्रान्तजीविनो विप्र दृश्यन्ते नु खगोत्तमाः}% ॥१९॥

\twolineshloka
{तस्मिंस्तथाविधे देशे कश्चिद्दैववशाद् वणिक्}
{निजसार्थपरिभ्रष्टः प्रविष्टो मरुजाङ्गले}% ॥२०॥

\twolineshloka
{बभ्रामोद्भ्रान्तहृदयः क्षुत्तृड्भ्यां श्रमपीडितः}
{क्व ग्रामः क्व जलं क्वाहं यास्यामि न बुबोध ह}% ॥२१॥

\twolineshloka
{अथ प्रेतान् ददर्शासौ क्षुत्तृषाव्याकुलेन्द्रियान्}
{उत्कटान् खलिनो भीमान् निर्मांसान् रौद्रदर्शनान्}% ॥२२॥

\twolineshloka
{प्रेतस्कन्धसमारूढम् एकं विकृतदर्शनम्}
{ददर्श बहुभिः प्रेतैः समन्तात् परिवारितम्}% ॥२३॥

\twolineshloka
{अगच्छमानमत्युग्रं प्रेतशब्दपुरःसरम्}
{प्रेतोऽपि दृष्ट्वा तां घोरामटवीमागतं नरम्}% ॥२४॥

\twolineshloka
{प्रेतस्कन्धान्महीं गत्वा तस्यान्तिकमुपागमत्}
{प्रणिपत्य वणिक्श्रेष्ठम् इदं वचनमब्रवीत्}% ॥२५॥

\twolineshloka
{अस्मिन् घोरतरे देशे प्रवेशो भवतः कथम्}
{तमुवाच वणिग्धीमान् सार्थभ्रष्टस्य मे वने}% ॥२६॥

\twolineshloka
{प्रवेशो दैवयोगेन पूर्वकर्मकृतेन च}
{तृषा मे बाधतेऽत्यर्थं क्षुधा चैव भृशं तथा}% ॥२७॥

\twolineshloka
{प्राणान्तिकमनुप्राप्तं वचनं भ्रमतीव मे}
{अत्रोपायं न पश्यामि जीवेयं येन केनचित्}% ॥२८॥

\twolineshloka
{इत्येवमुक्ते प्रेतस्तं वणिजं वाक्यमब्रवीत्}
{फुल्लां शमीं समाश्रित्य प्रतीक्ष त्वं मुहूर्तकम्}% ॥२९॥

\twolineshloka
{कृतातिथ्यो मया पश्चाद् गमिष्यसि यथासुखम्}
{एवमुक्तस्तथा चक्रे स वणिक् तृष्णयार्दितः}% ॥३०॥

\twolineshloka
{मध्याह्नसमये प्राप्ते प्रेतस्तं देशमागतः}
{फुल्लां सवृक्षां शीतोदां वारिधानीं मनोरमाम्}% ॥३१॥

\twolineshloka
{दध्योदनसमायुक्तां वर्द्धमानेन संयुताम्}
{अवतीर्य ततः स्वन्नं प्रादादतिथये तदा}% ॥३२॥

\twolineshloka
{स तत्राशनमात्रेण परं तृप्तित्वमागतः}
{वितृष्णो विज्वरश्चैव क्षणेन समपद्यत}% ॥३३॥

\twolineshloka
{ततश्च प्रेताः सम्प्राप्ताः सोऽस्माद्भागं क्रमाद्ददौ}
{दध्योदनात्सपानीयात् प्रेतास्तृप्तिं परां गताः}% ॥३४॥

\twolineshloka
{अतिथिं तर्पयित्वा तु प्रेतलोकं च सर्वतः}
{ततः स्वयं स बुभुजे भुक्तशेषं यथासुखम्}% ॥३५॥

\twolineshloka
{तस्य भुक्तवतः स्वन्नं पानीयं च क्षयं ययौ}
{प्रेताधिपं ततस्तं वै वणिग्वचनमब्रवीत्}% ॥३६॥

\twolineshloka
{आश्चर्यमेतत् परमं वनेऽस्मिन् प्रतिभाति मे}
{अन्नं पानं च परमं सम्प्राप्तं च कुतस्तव}% ॥३७॥

\twolineshloka
{स्वल्पेनैव तथाऽन्नेन त्वमेतांस्तु बहूनपि}
{अतर्पयः कथं त्वेते निर्मांसा भिन्नकुक्षयः}% ॥३८॥

\twolineshloka
{कथमस्यां सुघोरायामटव्यां च कृतालयाः}
{तदेतत् संशयं छिन्धि परं कौतूहलं मम}% ॥३९॥

\twolineshloka
{एवमुक्तः स वणिजा प्रेतो वचनमब्रवीत्}
{वाणिज्यसक्तस्य पुरा जन्मातीते ममानघ}% ॥४०॥

\twolineshloka
{सकले नगरे नास्ति ममान्यो हि दुरात्मकः}
{धनलोभान्न कस्यापि दत्ता भिक्षा मया तदा}% ॥४१॥

\twolineshloka
{सखा चैव ततश्चासीद् ब्राह्मणो गुणवान्मम}
{श्रावणद्वादशीयोगे मासि भाद्रपदे ततः}% ॥४२॥

\twolineshloka
{स कदाचिन्मया सार्द्धं तापीं नाम नदीं ययौ}
{तस्याश्च सङ्गमः पुण्यो यत्रासीच्चन्द्रभागया}% ॥४३॥

\twolineshloka
{चन्द्रभागा चन्द्रसुता तापी चैवार्कनन्दिनी}
{तयोः शीतोष्णसलिले प्रविवेश स ह द्विजः}% ॥४४॥

\twolineshloka
{श्रवणद्वादशीयोगे नराश्च समुपोषिताः}
{चन्द्रभागासुतो येन वारिधानीं ददुर्द्विजे}% ॥४५॥

\twolineshloka
{दध्योदनयुतां सार्द्धं सम्पूर्णैर्वर्द्धमानकैः}
{छत्रोपानद्युगं वस्त्रं प्रतिमां च तथा हरेः}% ॥४६॥

\twolineshloka
{प्रददौ विप्रमुख्येभ्यो हरस्याग्रे महामते}
{वित्तसंरक्षणार्थाय तस्यास्तीरे व्रते मया}% ॥४७॥

\twolineshloka
{सोपवासेन दत्तैका वारिधानी मनोरमा}
{तत्कृत्वाऽहं गृहं प्राप्तस्ततः कालेन केनचित्}% ॥४८॥

\twolineshloka
{पञ्चत्वमहमासाद्य नास्तिक्यात् प्रेततां गतः}
{अस्यामटव्यां घोरायां यथा ह्यहिकुलं तथा}% ॥४९॥

\twolineshloka
{श्रवणद्वादशीयोगे वारिधान्यर्पिता मया}
{सेयं मध्याह्नसमये लभ्यते च दिनेदिने}% ॥५०॥

\twolineshloka
{ब्रह्मस्वरूपिणः सर्वे पापाः प्रेतत्वमागताः}
{परदारगताः केचित् स्वामिद्रोहरताश्च ये}% ॥५१॥

\twolineshloka
{भूतप्रेतजरूपेण ते जाता ह्यत्र मानवाः}
{देशे मरुस्थले त्वस्मिन् ममैते मित्रतां गताः}% ॥५२॥

\twolineshloka
{अक्षयो भगवान् विष्णुः परमात्मा सनातनः}
{दीयते यत् समुद्दिश्य ह्यक्षयं तत् प्रकीर्तितम्}% ॥५३॥

\twolineshloka
{अक्षयेनापि चान्नेन तृप्ता एते पुनः पुनः}
{प्रेतत्वभावं दौर्बल्यं न विमुञ्चति कर्हिचित्}% ॥५४॥

\twolineshloka
{पूजयित्वाहमन्नैस्त्वामतिथिं समुपस्थितम्}
{प्रेतभावाद् विनिर्मुक्तो यास्यामि परमां गतिम्}% ॥५५॥

\twolineshloka
{मया विहीनाः किन्त्वेते वनेऽस्मिन् भृशदारुणे}
{पीडामनुभविष्यन्ति दारुणां कर्मयोनिजाम्}% ॥५६॥

\twolineshloka
{एतेषां तु महाभाग ममानुग्रहकाम्यया}
{प्रत्येकं नामगोत्राणि गृह्णीष्व लिखितानि च}% ॥५७॥

\twolineshloka
{अस्ति कक्षागता चैव तव सम्पुटिका शुभा}
{हिमवन्तमथासाद्य तत्र त्वं लप्स्यसे निधिम्}% ॥५८॥

\twolineshloka
{गयाशीर्षं ततो गत्वा श्राद्धं कुरु महामते}
{इत्याज्ञाप्य स वै प्रेतो वणिजं च यथासुखम्}% ॥५९॥

\twolineshloka
{विसर्जयामास तदा स वै प्रायात् समुत्सुकः}
{समासाद्य गृहं तत्र पश्चात् प्रायाद् हिमालयम्}% ॥६०॥

\twolineshloka
{ततो दृष्टं निधिं तत्र गृहीत्वा स समागतः}
{षष्ठांशं प्रतिगृह्याथ गयाशीर्षं ततोऽभ्यगात्}% ॥६१॥

\twolineshloka
{तत्र गत्वा गयायां स श्राद्धं कृत्वा महामतिः}
{प्रेतानां तु यथोद्दिष्टं श्राद्धं सम्यग्विधानतः}% ॥६२॥

\twolineshloka
{प्रत्येकं नामगोत्राणि गृहीत्वा पिण्डमुत्सृजत्}
{यस्य यस्य भवेत् श्राद्धं स करोति दिने वणिक्}% ॥६३॥

\twolineshloka
{स स तस्य तदा स्वप्ने दर्शयत्यात्मनस्तनुम्}
{ब्रवीति च महाभाग प्रसादाद् भवतोऽनघ}% ॥६४॥

\twolineshloka
{प्रेतभावं मया त्यक्तं प्राप्तोऽस्मि परमां गतिम्}
{एवं कृत्वा विधानेन गयाशीर्षं महामनाः}% ॥६५॥

\twolineshloka
{पश्चाज्जगाम स्वगृहं विष्णुं ध्यायन् पुनः पुनः}
{मासि भाद्रपदे प्राप्ते शुक्लपक्षे तथा सुधीः}% ॥६६॥

\twolineshloka
{श्रवणद्वादशीयोगे सङ्गमे सरितां पुनः}
{जगाम स महाबुद्धिः सर्वोपस्करसंयुतः}% ॥६७॥

\twolineshloka
{सङ्गमे सरितां स्नात्वा द्वादशीं तामुपोषितः}
{तत्र स्नात्वाऽपि दत्त्वा तु पूजयित्वा जनार्दनम्}% ॥६८॥

\twolineshloka
{अनन्तरं ब्राह्मणस्य ह्युपहारांस्तदा ददौ}
{शास्त्रोक्तेनापि विधिना ह्येकचित्तरतोऽपि सः}% ॥६९॥

\twolineshloka
{निवर्तयामास तदा वणिजो बुद्धिमान् स वै}
{वर्षे वर्षे तु सम्प्राप्ते मासि भाद्रपदे तथा}% ॥७०॥

\twolineshloka
{श्रवणद्वादशीयोगे सङ्गमे सरितां पुनः}
{एवं वै कृतवान् सर्वं विष्णुमुद्दिश्य सत्वरम्}% ॥७१॥

\twolineshloka
{कालेन चातिमहता पञ्चत्वं समुपागतः}
{अवाप परमं स्थानं दुर्लभं सर्वमानवैः}% ॥७२॥

\twolineshloka
{क्रीडतेऽद्यापि वैकुण्ठे विष्णुदूतैः स सेवितः}
{एवं कुरु त्वं भो ब्रह्मन् श्रवणद्वादशीव्रतम्}% ॥७३॥

\twolineshloka
{सर्वसौभाग्यदं चैव इह लोके परत्र च}
{सुबुद्धिजननं चैव सर्वपापहरं परम्}% ॥७४॥

\twolineshloka
{श्रवणद्वादशीयोगे यः कुर्याद् व्रतमीदृशम्}
{व्रतस्यास्य प्रभावेन विष्णुलोकं च गच्छति}% ॥७५॥

॥इति श्रीपाद्मे महापुराणे पञ्चपञ्चाशत्साहस्र्यां संहितायामुत्तरखण्डे उमापतिनारदसंवादे श्रवणद्वादशीव्रतं नाम सप्ततितमोऽध्यायः॥७०॥

\genMangalaShloka

\hyperref[sec:ekadashi_mahatmyam_padma_puranam]{\closesub}

\clearpage